\documentclass{ctexart} % 中文友好
\usepackage{amsmath}
\usepackage{lmodern}
\begin{document}
你好，\LaTeX！这是 XeLaTeX 自测。我是Jospeh.
这是一个引用范例\cite{lamport94}。
这是一个行内公式的测试：$ \frac{z}{d} $
\begin{equation}\label{1}
    \int_{0}^{1}f(t)dt = \iint_{D}g(x,y)dxdy.
\end{equation}
\begin{equation}\label{2}
\begin{split}
\frac{1}{2} (\sin(x+y) + \sin(x-y)) =& \frac{1}{2}(\sin x \cos y + \cos x \sin y)\\
& + \frac{1}{2} (\sin x \cos y - \cos x \sin y)\\
=& \sin x \cos y.
\end{split}
\end{equation}
麦克斯韦方程组：
麦克斯韦方程组是经典电磁学的核心，它用四个偏微分方程统一描述了电场、磁场与电荷、电流之间的相互关系，并预言了电磁波的存在。
\begin{equation}
\begin{split}
    \nabla \cdot \mathbf{D} &= \rho, \\
    \nabla \cdot \mathbf{B} &= 0, \\
    \nabla \times \mathbf{E} &= -\frac{\partial \mathbf{B}}{\partial t}, \\
    \nabla \times \mathbf{H} &= \mathbf{J} + \frac{\partial \mathbf{D}}{\partial t}.
\end{split}
\label{maxwell}
\end{equation}
\subsection*{方程组解释}
\begin{itemize}
    \item \textbf{高斯定律}（式一）：$\nabla \cdot \mathbf{D} = \rho$。电场的有源性质——电荷是电场的源，穿过闭合曲面的电通量正比于曲面内所含的自由电荷量。
    \item \textbf{高斯磁定律}（式二）：$\nabla \cdot \mathbf{B} = 0$。磁场的无源性质——不存在磁单极子，磁力线总是闭合的，穿过任意闭合曲面的磁通量恒为零。
    \item \textbf{法拉第电磁感应定律}（式三）：$\nabla \times \mathbf{E} = -\frac{\partial \mathbf{B}}{\partial t}$。变化的磁场产生涡旋电场——这是发电机、变压器工作的基本原理。
    \item \textbf{安培-麦克斯韦定律}（式四）：$\nabla \times \mathbf{H} = \mathbf{J} + \frac{\partial \mathbf{D}}{\partial t}$。电流和变化的电场都能产生磁场，其中位移电流项 $\frac{\partial \mathbf{D}}{\partial t}$ 是麦克斯韦的关键贡献，它使得电磁波的存在成为可能。
\end{itemize}
\subsection*{主要应用}
\begin{enumerate}
    \item \textbf{电磁波与通信}：麦克斯韦预言了电磁波的存在，后被赫兹实验证实。无线电、微波、光通信（光纤）、雷达均基于此。
    \item \textbf{电力系统}：发电机（法拉第定律）与变压器将机械能转化为电能并实现电压变换，构成现代电网的基础。
    \item \textbf{电磁兼容（EMC）}：电子设备设计中需控制电磁辐射与抗干扰能力，确保各设备在共用电磁环境中正常工作。
    \item \textbf{医用成像}：MRI（磁共振成像）利用强磁场与射频脉冲，基于核磁共振原理获取人体内部结构图像。
    \item \textbf{天线设计}：从手机天线到射电望远镜，天线辐射方向图与阻抗匹配均由麦克斯韦方程组导出。
    \item \textbf{光学与光电子}：可见光本质上是特定频段的电磁波，反射、折射、衍射、偏振等现象均可由麦克斯韦方程组统一描述。
\end{enumerate}
\subsection*{重要推导与推论}
从麦克斯韦方程组出发，可以推导出一系列具有深远影响的结论：
\begin{enumerate}
    \item \textbf{电磁波方程}。在无源区域（$\rho = 0$，$\mathbf{J} = 0$）中，对法拉第定律取旋度并代入安培-麦克斯韦定律，得到波动方程：
    \begin{equation}
        \nabla^2 \mathbf{E} - \mu\varepsilon \frac{\partial^2 \mathbf{E}}{\partial t^2} = 0,
    \end{equation}
    其传播速度 $v = 1/\sqrt{\mu\varepsilon}$。在真空中即光速 $c = 1/\sqrt{\mu_0\varepsilon_0}$，从而揭示光就是电磁波，实现了电学、磁学与光学的统一。
    \item \textbf{电荷守恒定律（连续性方程）}。对安培-麦克斯韦定律取散度，利用矢量恒等式和电荷守恒，得到：
    \begin{equation}
        \nabla \cdot \mathbf{J} + \frac{\partial \rho}{\partial t} = 0.
    \end{equation}
    这一方程表明电荷既不能创生也不能消灭，仅能从一处流向另一处。
    \item \textbf{坡印廷定理（能量守恒）}。将法拉第定律与安培-麦克斯韦定律联立，可以导出电磁场的能量守恒关系：
    \begin{equation}
        \frac{\partial u}{\partial t} + \nabla \cdot \mathbf{S} = -\mathbf{J} \cdot \mathbf{E},
    \end{equation}
    其中能量密度 $u = \frac{1}{2}(\mathbf{E} \cdot \mathbf{D} + \mathbf{B} \cdot \mathbf{H})$，坡印廷矢量 $\mathbf{S} = \mathbf{E} \times \mathbf{H}$ 表示电磁能的流动方向和通量密度。
    \item \textbf{库仑定律的等价性}。在静电场条件下（$\partial/\partial t = 0$），高斯定律 $\nabla \cdot \mathbf{E} = \rho/\varepsilon_0$ 通过散度定理等价于库仑定律：
    \begin{equation}
        \mathbf{F} = \frac{1}{4\pi\varepsilon_0}\frac{q_1 q_2}{r^2}\hat{\mathbf{r}}.
    \end{equation}
    \item \textbf{边界条件}。在两种不同介质的交界面处，麦克斯韦方程组的积分形式导出以下边界条件：
    \begin{align}
        D_{1n} - D_{2n} &= \sigma_f, & B_{1n} - B_{2n} &= 0,\\
        E_{1t} - E_{2t} &= 0,      & H_{1t} - H_{2t} &= K_f.
    \end{align}
    这些条件在电磁波反射与折射、波导和谐振腔分析中至关重要。
    \item \textbf{电磁势与规范不变性}。由 $\nabla \cdot \mathbf{B} = 0$ 引入矢量势 $\mathbf{A}$ 使得 $\mathbf{B} = \nabla \times \mathbf{A}$，代入法拉第定律得 $\mathbf{E} = -\nabla\varphi - \partial\mathbf{A}/\partial t$。规范变换：
    \begin{equation}
        \mathbf{A} \to \mathbf{A} + \nabla\Lambda,\quad \varphi \to \varphi - \frac{\partial\Lambda}{\partial t},
    \end{equation}
    保持电磁场不变，这一对称性在量子电动力学（QED）中起着基石作用。
\end{enumerate}
%TODO 测试注释效果
\clearpage
\bibliographystyle{plain}
\bibliography{refs}
\end{document}
